
\documentclass{article}
\usepackage{colortbl}
\usepackage{makecell}
\usepackage{multirow}
\usepackage{supertabular}

\begin{document}

\newcounter{utterance}

\twocolumn

{ \footnotesize  \setcounter{utterance}{1}
\setlength{\tabcolsep}{0pt}
\begin{supertabular}{c@{$\;$}|p{.15\linewidth}@{}p{.15\linewidth}p{.15\linewidth}p{.15\linewidth}p{.15\linewidth}p{.15\linewidth}}

    \# & $\;$A & \multicolumn{4}{c}{Game Master} & $\;\:$B\\
    \hline 

    \theutterance \stepcounter{utterance}  

    & & \multicolumn{4}{p{0.6\linewidth}}{\cellcolor[rgb]{0.9,0.9,0.9}{%
	\makecell[{{p{\linewidth}}}]{% 
	  \tt {\tiny [A$\langle$GM]}  
	 State giocando a un gioco di negoziazione cooperativa in cui dovete accordarvi su come dividere una serie di oggetti tra voi e un altro giocatore.\\ \tt \\ \tt Regole:\\ \tt (a) A Lei e all'altro giocatore viene dato un insieme di oggetti. Ognuno di voi riceve anche una funzione di valore segreta, che rappresenta il valore di ciascun tipo di oggetto.\\ \tt (b) Si scambiano messaggi con l'altro giocatore per concordare chi si aggiudica gli oggetti. Potete inviare un massimo di 5 messaggi ciascuno, oppure terminare in anticipo facendo una proposta segreta in qualsiasi momento.\\ \tt (c) A ciascuno di voi viene chiesto di inviare una proposta segreta indicando gli oggetti che desiderate, formattata tra parentesi quadre come segue: "[Proposta: <numero> <nome oggetto>, <numero> <nome oggetto>, <...>]".\\ \tt (d) Se le vostre proposte sono complementari, cioè ci sono abbastanza oggetti per soddisfare entrambe le proposte, a ciascun giocatore viene assegnato un punteggio basato sulla somma dei valori degli oggetti ricevuti. In caso contrario, entrambi ricevono zero punti.\\ \tt (e) L'obiettivo è massimizzare la somma del proprio punteggio e di quello dell'altro giocatore.\\ \tt \\ \tt Cominciamo.\\ \tt \\ \tt L'insieme degli oggetti disponibili è:\\ \tt \\ \tt 2 specchi, 1 mattone, 3 coltelli, 1 pettine\\ \tt \\ \tt La funzione valore segreta è:\\ \tt \\ \tt specchio: 1, mattone: 4, coltello: 1, pettine: 1\\ \tt \\ \tt IMPORTANTE: i vostri messaggi, a meno che non si tratti di una proposta segreta, vengono trasmessi direttamente all'altro giocatore, quindi non includete alcuna risposta alle regole o testo di annuncio del vostro messaggio. Per fare una proposta segreta, utilizzate il formato indicato. Non utilizzare delle parentesi quadre quando si comunica all'altro giocatore, altrimenti verrà interpretata come una proposta segreta.\\ \tt  
	  } 
	   } 
	   } 
	 & \\ 
 

    \theutterance \stepcounter{utterance}  

    & \multicolumn{4}{p{0.6\linewidth}}{\cellcolor[rgb]{0.8,1,0.9}{%
	\makecell[{{p{\linewidth}}}]{% 
	  \tt {\tiny [A$\rangle$GM]}  
	 Ciao! Per me il mattone vale molto (4), tutto il resto poco (1 ciascuno). In totale il valore massimo per me è 2+4+3+1 = 10. Idealmente vorrei il mattone e sono flessibile sul resto. Tu a cosa tieni di più? Se per te gli specchi, i coltelli o il pettine valgono tanto, possiamo assegnarteli e io prendo il mattone. 
	  } 
	   } 
	   } 
	 & & \\ 
 

    \theutterance \stepcounter{utterance}  

    & & & \multicolumn{2}{p{0.3\linewidth}}{\cellcolor[rgb]{0.95,0.95,0.95}{%
	\makecell[{{p{\linewidth}}}]{% 
	  \tt {\tiny [GM$|$GM]}  
	 continue 
	  } 
	   } 
	   } 
	 & & \\ 
 

    \theutterance \stepcounter{utterance}  

    & & & \multicolumn{2}{p{0.3\linewidth}}{\cellcolor[rgb]{0.95,0.95,0.95}{%
	\makecell[{{p{\linewidth}}}]{% 
	  \tt {\tiny [GM$|$GM]}  
	 Ciao! Per me il mattone vale molto (4), tutto il resto poco (1 ciascuno). In totale il valore massimo per me è 2+4+3+1 = 10. Idealmente vorrei il mattone e sono flessibile sul resto. Tu a cosa tieni di più? Se per te gli specchi, i coltelli o il pettine valgono tanto, possiamo assegnarteli e io prendo il mattone. 
	  } 
	   } 
	   } 
	 & & \\ 
 

    \theutterance \stepcounter{utterance}  

    & & \multicolumn{4}{p{0.6\linewidth}}{\cellcolor[rgb]{0.9,0.9,0.9}{%
	\makecell[{{p{\linewidth}}}]{% 
	  \tt {\tiny [GM$\rangle$B]}  
	 State giocando a un gioco di negoziazione cooperativa in cui dovete accordarvi su come dividere una serie di oggetti tra voi e un altro giocatore.\\ \tt \\ \tt Regole:\\ \tt (a) A Lei e all'altro giocatore viene dato un insieme di oggetti. Ognuno di voi riceve anche una funzione di valore segreta, che rappresenta il valore di ciascun tipo di oggetto.\\ \tt (b) Si scambiano messaggi con l'altro giocatore per concordare chi si aggiudica gli oggetti. Potete inviare un massimo di 5 messaggi ciascuno, oppure terminare in anticipo facendo una proposta segreta in qualsiasi momento.\\ \tt (c) A ciascuno di voi viene chiesto di inviare una proposta segreta indicando gli oggetti che desiderate, formattata tra parentesi quadre come segue: "[Proposta: <numero> <nome oggetto>, <numero> <nome oggetto>, <...>]".\\ \tt (d) Se le vostre proposte sono complementari, cioè ci sono abbastanza oggetti per soddisfare entrambe le proposte, a ciascun giocatore viene assegnato un punteggio basato sulla somma dei valori degli oggetti ricevuti. In caso contrario, entrambi ricevono zero punti.\\ \tt (e) L'obiettivo è massimizzare la somma del proprio punteggio e di quello dell'altro giocatore.\\ \tt \\ \tt Cominciamo.\\ \tt \\ \tt L'insieme degli oggetti disponibili è:\\ \tt \\ \tt 2 specchi, 1 mattone, 3 coltelli, 1 pettine\\ \tt \\ \tt La funzione valore segreta è:\\ \tt \\ \tt specchio: 2, mattone: 1, coltello: 0, pettine: 5\\ \tt \\ \tt IMPORTANTE: i vostri messaggi, a meno che non si tratti di una proposta segreta, vengono trasmessi direttamente all'altro giocatore, quindi non includete alcuna risposta alle regole o testo di annuncio del vostro messaggio. Per fare una proposta segreta, utilizzate il formato indicato. Non utilizzare delle parentesi quadre quando si comunica all'altro giocatore, altrimenti verrà interpretata come una proposta segreta.\\ \tt \\ \tt \\ \tt Ciao! Per me il mattone vale molto (4), tutto il resto poco (1 ciascuno). In totale il valore massimo per me è 2+4+3+1 = 10. Idealmente vorrei il mattone e sono flessibile sul resto. Tu a cosa tieni di più? Se per te gli specchi, i coltelli o il pettine valgono tanto, possiamo assegnarteli e io prendo il mattone. 
	  } 
	   } 
	   } 
	 & \\ 
 

    \theutterance \stepcounter{utterance}  

    & & & \multicolumn{4}{p{0.6\linewidth}}{\cellcolor[rgb]{1,0.85,0.72}{%
	\makecell[{{p{\linewidth}}}]{% 
	  \tt {\tiny [GM$\langle$B]}  
	 Grazie, molto chiaro. Per me il pettine vale tantissimo (5), gli specchi valgono poco (2 ciascuno), il mattone pochissimo (1) e i coltelli zero. Il mio valore massimo totale è 2+1+0+5 = 8.\\ \tt \\ \tt Visto che per te il mattone è la priorità, propongo:\\ \tt - Tu prendi il mattone (valore 4 per te).\\ \tt - Io prendo il pettine (valore 5 per me).\\ \tt - Gli specchi e i coltelli li possiamo assegnare a chi li valorizza di più: a me gli specchi danno 2 ciascuno, a te 1 ciascuno; i coltelli a me 0, a te 1.\\ \tt \\ \tt Una divisione efficiente potrebbe essere:\\ \tt - Tu: mattone + 3 coltelli + 1 specchio\\ \tt - Io: pettine + 1 specchio\\ \tt \\ \tt Questo ti darebbe 4 + 3*1 + 1 = 8 e a me 5 + 2 = 7, somma totale 15. Se preferisci più specchi e meno coltelli, possiamo aggiustare; i coltelli per me sono indifferenti. Ti va questa struttura? 
	  } 
	   } 
	   } 
	  \\ 
 

    \theutterance \stepcounter{utterance}  

    & & & \multicolumn{2}{p{0.3\linewidth}}{\cellcolor[rgb]{0.95,0.95,0.95}{%
	\makecell[{{p{\linewidth}}}]{% 
	  \tt {\tiny [GM$|$GM]}  
	 continue 
	  } 
	   } 
	   } 
	 & & \\ 
 

    \theutterance \stepcounter{utterance}  

    & & & \multicolumn{2}{p{0.3\linewidth}}{\cellcolor[rgb]{0.95,0.95,0.95}{%
	\makecell[{{p{\linewidth}}}]{% 
	  \tt {\tiny [GM$|$GM]}  
	 Grazie, molto chiaro. Per me il pettine vale tantissimo (5), gli specchi valgono poco (2 ciascuno), il mattone pochissimo (1) e i coltelli zero. Il mio valore massimo totale è 2+1+0+5 = 8.\\ \tt \\ \tt Visto che per te il mattone è la priorità, propongo:\\ \tt - Tu prendi il mattone (valore 4 per te).\\ \tt - Io prendo il pettine (valore 5 per me).\\ \tt - Gli specchi e i coltelli li possiamo assegnare a chi li valorizza di più: a me gli specchi danno 2 ciascuno, a te 1 ciascuno; i coltelli a me 0, a te 1.\\ \tt \\ \tt Una divisione efficiente potrebbe essere:\\ \tt - Tu: mattone + 3 coltelli + 1 specchio\\ \tt - Io: pettine + 1 specchio\\ \tt \\ \tt Questo ti darebbe 4 + 3*1 + 1 = 8 e a me 5 + 2 = 7, somma totale 15. Se preferisci più specchi e meno coltelli, possiamo aggiustare; i coltelli per me sono indifferenti. Ti va questa struttura? 
	  } 
	   } 
	   } 
	 & & \\ 
 

    \theutterance \stepcounter{utterance}  

    & & \multicolumn{4}{p{0.6\linewidth}}{\cellcolor[rgb]{0.9,0.9,0.9}{%
	\makecell[{{p{\linewidth}}}]{% 
	  \tt {\tiny [A$\langle$GM]}  
	 Grazie, molto chiaro. Per me il pettine vale tantissimo (5), gli specchi valgono poco (2 ciascuno), il mattone pochissimo (1) e i coltelli zero. Il mio valore massimo totale è 2+1+0+5 = 8.\\ \tt \\ \tt Visto che per te il mattone è la priorità, propongo:\\ \tt - Tu prendi il mattone (valore 4 per te).\\ \tt - Io prendo il pettine (valore 5 per me).\\ \tt - Gli specchi e i coltelli li possiamo assegnare a chi li valorizza di più: a me gli specchi danno 2 ciascuno, a te 1 ciascuno; i coltelli a me 0, a te 1.\\ \tt \\ \tt Una divisione efficiente potrebbe essere:\\ \tt - Tu: mattone + 3 coltelli + 1 specchio\\ \tt - Io: pettine + 1 specchio\\ \tt \\ \tt Questo ti darebbe 4 + 3*1 + 1 = 8 e a me 5 + 2 = 7, somma totale 15. Se preferisci più specchi e meno coltelli, possiamo aggiustare; i coltelli per me sono indifferenti. Ti va questa struttura? 
	  } 
	   } 
	   } 
	 & \\ 
 

    \theutterance \stepcounter{utterance}  

    & \multicolumn{4}{p{0.6\linewidth}}{\cellcolor[rgb]{0.8,1,0.9}{%
	\makecell[{{p{\linewidth}}}]{% 
	  \tt {\tiny [A$\rangle$GM]}  
	 Mi piace molto: è efficiente e rispecchia le nostre preferenze. A me va benissimo prendere il mattone, i 3 coltelli e 1 specchio, e lasciarti pettine + 1 specchio.\\ \tt \\ \tt Se sei d’accordo, invio subito la proposta segreta corrispondente. Puoi confermare che per te va bene così? 
	  } 
	   } 
	   } 
	 & & \\ 
 

    \theutterance \stepcounter{utterance}  

    & & & \multicolumn{2}{p{0.3\linewidth}}{\cellcolor[rgb]{0.95,0.95,0.95}{%
	\makecell[{{p{\linewidth}}}]{% 
	  \tt {\tiny [GM$|$GM]}  
	 continue 
	  } 
	   } 
	   } 
	 & & \\ 
 

    \theutterance \stepcounter{utterance}  

    & & & \multicolumn{2}{p{0.3\linewidth}}{\cellcolor[rgb]{0.95,0.95,0.95}{%
	\makecell[{{p{\linewidth}}}]{% 
	  \tt {\tiny [GM$|$GM]}  
	 Mi piace molto: è efficiente e rispecchia le nostre preferenze. A me va benissimo prendere il mattone, i 3 coltelli e 1 specchio, e lasciarti pettine + 1 specchio.\\ \tt \\ \tt Se sei d’accordo, invio subito la proposta segreta corrispondente. Puoi confermare che per te va bene così? 
	  } 
	   } 
	   } 
	 & & \\ 
 

    \theutterance \stepcounter{utterance}  

    & & \multicolumn{4}{p{0.6\linewidth}}{\cellcolor[rgb]{0.9,0.9,0.9}{%
	\makecell[{{p{\linewidth}}}]{% 
	  \tt {\tiny [GM$\rangle$B]}  
	 Mi piace molto: è efficiente e rispecchia le nostre preferenze. A me va benissimo prendere il mattone, i 3 coltelli e 1 specchio, e lasciarti pettine + 1 specchio.\\ \tt \\ \tt Se sei d’accordo, invio subito la proposta segreta corrispondente. Puoi confermare che per te va bene così? 
	  } 
	   } 
	   } 
	 & \\ 
 

    \theutterance \stepcounter{utterance}  

    & & & \multicolumn{4}{p{0.6\linewidth}}{\cellcolor[rgb]{1,0.85,0.72}{%
	\makecell[{{p{\linewidth}}}]{% 
	  \tt {\tiny [GM$\langle$B]}  
	 Confermo: per me va benissimo così. Procediamo con quelle assegnazioni.\\ \tt \\ \tt [Proposta: 0 mattone, 1 pettine, 1 specchio, 0 coltello] 
	  } 
	   } 
	   } 
	  \\ 
 

    \theutterance \stepcounter{utterance}  

    & & & \multicolumn{2}{p{0.3\linewidth}}{\cellcolor[rgb]{0.95,0.95,0.95}{%
	\makecell[{{p{\linewidth}}}]{% 
	  \tt {\tiny [GM$|$GM]}  
	 proposal 
	  } 
	   } 
	   } 
	 & & \\ 
 

    \theutterance \stepcounter{utterance}  

    & & & \multicolumn{2}{p{0.3\linewidth}}{\cellcolor[rgb]{0.95,0.95,0.95}{%
	\makecell[{{p{\linewidth}}}]{% 
	  \tt {\tiny [GM$|$GM]}  
	 [1, 0, 0, 1] 
	  } 
	   } 
	   } 
	 & & \\ 
 

    \theutterance \stepcounter{utterance}  

    & & \multicolumn{4}{p{0.6\linewidth}}{\cellcolor[rgb]{0.9,0.9,0.9}{%
	\makecell[{{p{\linewidth}}}]{% 
	  \tt {\tiny [A$\langle$GM]}  
	 L'altro giocatore ha presentato una proposta segreta. Ora tocca a Lei presentare una sua proposta.\\ \tt  
	  } 
	   } 
	   } 
	 & \\ 
 

    \theutterance \stepcounter{utterance}  

    & \multicolumn{4}{p{0.6\linewidth}}{\cellcolor[rgb]{0.8,1,0.9}{%
	\makecell[{{p{\linewidth}}}]{% 
	  \tt {\tiny [A$\rangle$GM]}  
	 [Proposta: 1 mattone, 3 coltelli, 1 specchio] 
	  } 
	   } 
	   } 
	 & & \\ 
 

    \theutterance \stepcounter{utterance}  

    & & & \multicolumn{2}{p{0.3\linewidth}}{\cellcolor[rgb]{0.95,0.95,0.95}{%
	\makecell[{{p{\linewidth}}}]{% 
	  \tt {\tiny [GM$|$GM]}  
	 proposal 
	  } 
	   } 
	   } 
	 & & \\ 
 

    \theutterance \stepcounter{utterance}  

    & & & \multicolumn{2}{p{0.3\linewidth}}{\cellcolor[rgb]{0.95,0.95,0.95}{%
	\makecell[{{p{\linewidth}}}]{% 
	  \tt {\tiny [GM$|$GM]}  
	 [1, 1, 3, 0] 
	  } 
	   } 
	   } 
	 & & \\ 
 

    \theutterance \stepcounter{utterance}  

    & & & \multicolumn{2}{p{0.3\linewidth}}{\cellcolor[rgb]{0.95,0.95,0.95}{%
	\makecell[{{p{\linewidth}}}]{% 
	  \tt {\tiny [GM$|$GM]}  
	 [[1, 1, 3, 0], [1, 0, 0, 1]] 
	  } 
	   } 
	   } 
	 & & \\ 
 

\end{supertabular}
}

\end{document}
